% SoftwareX Original Software Publication, based on the official 2026 template.
% Evidence fields are populated by ../build_paper.py, not manually estimated.
\documentclass[preprint,12pt,a4paper]{elsarticle}
\usepackage[T1]{fontenc}
\usepackage{amssymb,amsmath,booktabs,array,hyperref,url}
\expandafter\def\expandafter\UrlBreaks\expandafter{\UrlBreaks\do-}
\setlength{\parindent}{0pt}
\setlength{\emergencystretch}{2em}
\hypersetup{pdftitle={ZeroRun: Reproducible test-result reuse for AI coding tools},pdfauthor={Jishan Kapoor},colorlinks=false,pdfborder={0 0 0}}
\journal{SoftwareX}
\newcommand{\code}[1]{\texttt{#1}}
\begin{document}
\begin{frontmatter}
\title{ZeroRun: Reproducible test-result reuse for AI coding tools}
\author{Jishan Kapoor\corref{cor1}}
\ead{kapoorjishan2@gmail.com}
\cortext[cor1]{Corresponding author.}
\address{Independent researcher, Toronto, Canada}
\begin{abstract}
ZeroRun is open-source software for reusing the identified successful status of a deterministic test command in AI coding-tool integrations. It combines declared input and runtime identities, authenticated results, isolated fresh execution, and command-line and coding-tool interfaces. A read-only hit path avoids source staging. The release includes fresh-result comparisons, an independent input-boundary oracle, a descriptive audit of public agent trajectories, and runnable client examples. The four-target study, completed through an explicit post-crash recovery, reconciles 168 fresh comparisons and 96 reused successes. On one synthetic fixture, a controlled Codex lifecycle distinguished reused status from fresh validation. The software enables investigation of validation cost and reuse boundaries; it does not replay cached output or establish autonomous-agent acceleration.
\end{abstract}
\begin{keyword}
AI coding tools \sep test-result reuse \sep research software \sep reproducibility \sep software testing
\end{keyword}
\end{frontmatter}
\clearpage

\section*{Required Metadata}
\section*{Current code version}
\begin{table}[!ht]
\small
\begin{tabular}{|l|>{\raggedright\arraybackslash}p{5.8cm}|>{\raggedright\arraybackslash}p{5.8cm}|}
\hline
\textbf{Nr.} & \textbf{Code metadata description} & \textbf{Value} \\
\hline
C1 & Current code version & ZeroRun 0.5.1; research release 2026-09-07 \\
\hline
C2 & Permanent GitHub link to code/repository used for this code version & \href{https://github.com/floxy-21/zerorun-research/tree/bdf7fdc1309381b1c6762b2189547d1a1da87263}{\texttt{floxy-21/zerorun-research}}; linked immutable source/evidence snapshot \\
\hline
C3 & Legal Code License & MIT (ZeroRun code); third-party licenses retained; trajectory data CC-BY-4.0 \\
\hline
C4 & Code versioning system used & Git \\
\hline
C5 & Software code languages, tools, and services used & Python; pytest; Docker; CLI; Model Context Protocol (MCP); Codex skill \\
\hline
C6 & Compilation requirements, operating environments \& dependencies & Runtime: Python $\geq$3.10; no core Python dependencies. Whole-task reuse: Linux/amd64 and Docker. Canonical analysis: CPython 3.12--3.14. Experiments: digest-pinned Python 3.12 container and frozen dependencies. \\
\hline
C7 & If available Link to developer documentation/manual & \href{https://github.com/floxy-21/zerorun-research/blob/bdf7fdc1309381b1c6762b2189547d1a1da87263/README.md}{Version-pinned README: installation, interfaces, reproduction, and limitations} \\
\hline
C8 & Support email for questions & \href{mailto:kapoorjishan2@gmail.com}{kapoorjishan2@gmail.com} \\
\hline
\end{tabular}
\caption{Code metadata. The public release preserves the tested runtime source and provides a separate, documented research harness.}
\end{table}
\clearpage

\section{Motivation and significance}
Repository-level AI coding tools use test results when deciding whether an edit is satisfactory. SWE-bench and SWE-agent make this interaction a concrete research setting \cite{jimenez2024swebench,yang2024sweagent}. A reused result must be distinguishable from a new execution: a previous success is not a fresh transcript, and the same shell command can observe different source, dependencies, or external state.

ZeroRun exposes the previous successful exit status of a deterministic test command under an unchanged declared execution identity. Its users are testing-tool integrators and researchers evaluating whether avoided execution exceeds validation cost. A coding tool revisiting a previously tested source state can request its status, while requesting fresh execution for new diagnostics. ZeroRun neither learns a caching policy nor predicts patch correctness.

Regression-test selection avoids tests considered unaffected by a change \cite{rothermel1997safe,yoo2012regression}. Existing build and LLM-tool caches already address result reuse (Table~\ref{tab:contracts}); eidos provides adjacent LLM function-integration software \cite{aldanamartin2025eidos}. Neither content hashing nor attaching a cache to an AI interface is a new algorithm here.

The contribution combines a repository-bound status interface, a hit path that avoids source staging, and an executable evaluation kit linking input boundaries, fresh outcomes, and full request cost. Its distinguishing unit is an operator-qualified, result-only test task, not an arbitrary tool result or sandbox history. Integrators can separately inspect three questions: does earlier status agree with fresh execution, is reuse eligible, and is it worthwhile? The software supplies those decision and measurement components without requiring a claim of universal acceleration.

\begin{table}[htbp]
\centering\small
\begin{tabular}{@{}>{\raggedright\arraybackslash}p{0.17\linewidth}>{\raggedright\arraybackslash}p{0.35\linewidth}>{\raggedright\arraybackslash}p{0.38\linewidth}@{}}
\toprule
System & Reuse basis & Returned information / execution state \\
\midrule
Bazel \cite{bazelremotecache} & Build-action hashes with declared inputs and outputs & Action metadata, output files, stdout and stderr. \\
ToolCaching \cite{zhai2026toolcaching} & Tool name and serialized arguments; feature/TTL-based admission & Tool results and metadata; requests classified as COMMAND excluded. \\
TVCache \cite{kumar2026tvcache} & Tool-call prefixes; optional state-preserving annotations & Tool values; selective sandbox snapshots support continued rollout execution. \\
ZeroRun & Reviewed local test inputs, command and runtime identity & Identified earlier success; no historical transcript or restored output artifacts. \\
\bottomrule
\end{tabular}
\caption{Operating contracts, not a performance ranking. ZeroRun evaluates a narrower result-only task contract; no head-to-head superiority is established.}
\label{tab:contracts}
\end{table}

\section{Software description}
\subsection{Architecture and execution contract}
Figure~\ref{fig:architecture} shows the package's manifest validation, input fingerprinting, result storage, isolated execution, and interfaces. A version-2 whole-task manifest identifies a command, declared input paths, environment policy, and digest-pinned Linux/amd64 container. Its supported cache contract requires a result-only computation, no output artifacts, and no cached stdout or stderr. The input closure and determinism are operator assumptions, not facts inferred from passing tests.

\begin{figure}[htbp]
\centering
\setlength{\fboxsep}{7pt}
\fbox{\parbox{0.88\linewidth}{\centering
CLI / Codex skill / configured-task MCP request\\[4pt]
$\downarrow$\\[4pt]
Interface-specific checks; runtime inspection; declared-input fingerprint\\[4pt]
$\downarrow$\\[4pt]
Task-key lock and authenticated result lookup\\[6pt]
\begin{tabular}{p{0.41\linewidth}p{0.41\linewidth}}
\centering Eligible hit & \centering Miss or explicit verification\tabularnewline
\centering Recompute fingerprint; require identity agreement & \centering Construct private source snapshot; execute isolated command\tabularnewline
\centering Return identified reused success & \centering Recheck source; return fresh result; publish only eligible success\tabularnewline
\end{tabular}\\[6pt]
Unsupported conditions produce fresh-required or refusal responses; no silent success.
}}
\caption{Whole-file result-reuse path. A hit returns status metadata, not the prior transcript. This explanatory diagram was revised as editable LaTeX with OpenAI Codex desktop 26.901.6511.0 assistance and checked against the described implementation; it is not experimental image data.}
\label{fig:architecture}
\end{figure}

The identity includes declared content, relevant command-source content, directory modes, command and environment policy, and inspected runtime identity. A matching cache filename is insufficient: the entry must authenticate and satisfy the configured contract. On a miss, ZeroRun stages a private source copy and runs it read-only in a resource-constrained, network-disabled container. A post-execution check detects source drift before publishing a success. Failed commands do not become reusable successes. Explicit verification can quarantine a stored result that conflicts with fresh execution.

The optimized whole-file hit path computes the full fingerprint, locks and authenticates the entry, then independently recomputes the full fingerprint. It returns a hit only if both identities and their payloads agree. No repository code executes and no execution copy is constructed on this path. Misses, forced execution, verification, and symbol-projection tasks retain the existing snapshot path. The optimization changes the order of work, not the cache eligibility threshold.

This is conditional correctness, not a soundness proof for arbitrary Python. An undeclared dependency, nondeterministic test, or violated host assumption can invalidate reuse. Two filesystem scans are not an atomic snapshot against arbitrary concurrent host mutations. Unsupported or insufficiently reviewed work must execute fresh or remain refused.

\subsection{Interfaces and functionality}
The CLI exposes structured results identifying reused success, fresh success, fresh failure, uncertainty, and refusal. The Codex skill and repository-bound MCP server expose the same distinction. Normal agent-facing reuse additionally requires exact-configuration authorization outside repository-controlled files; generating observation candidates does not authorize reuse. Laboratory fixtures used by the experiments authenticate cache records but are not production operator reviews, and their private keys are excluded from the artifact.

The package also contains a fine-grained pytest path. Because omitted tests can affect fixtures, callbacks, shared state, or ordering, it conservatively requires fresh execution when qualification is incomplete. That path is included for reproducibility, but it is not the demonstrated acceleration mechanism in this article. Whole-task reuse does not require individual tests to be independent; it requires the complete requested command to satisfy the stated contract.

Hits do not reconstruct stdout or stderr. Clients requiring new diagnostics, coverage, or generated files must execute fresh. The public operating guide and manifest reference document installation, operator review, authorization, and this boundary; automated interface checks do not establish external adoption.

\section{Illustrative examples}
\subsection{Repeat, failure, and restoration}
The executable example starts with an empty store and a fixed pytest target. It executes a seed, reuses unchanged successes, executes an added failure twice, then finds the earlier success after byte restoration. The seven-response sequence is:
\begin{center}
\small
\begin{tabular}{ll}
\toprule
Request & Expected whole-task response \\
\midrule
Seed & \code{MISS\_EXECUTED} \\
Two unchanged repeats & \code{HIT\_REUSED} each \\
Added failure; repeated failure & \code{MISS\_FAILED} each \\
Restoration; restored repeat & \code{HIT\_REUSED} each \\
\bottomrule
\end{tabular}
\end{center}
The example records an independent fresh full-target check after every request, including hits. It is a constructed validation sequence, not an estimate of the repetition frequency in real development.

For a configured target named \code{tests}, the CLI form is:
\begin{verbatim}
zerorun --manifest .zerorun.json --json explain tests
zerorun --manifest .zerorun.json --json run tests
zerorun --manifest .zerorun.json --json run tests --verify
\end{verbatim}
These commands inspect eligibility, request execution/reuse, and verify a stored result through fresh execution, respectively. The target must already exist in a reviewed manifest; the commands neither establish input completeness nor override authority. The harness uses laboratory-only configuration and retains the complete result JSON.

\subsection{Client decisions and integration evidence}
The reference consumer checks response consistency, \code{status}, \code{exit\_code}, and tool/protocol errors: \code{mode=reuse} alone is not success. It labels previous success explicitly, routes fresh failures to diagnostics, and requests fresh execution when required; it does not authorize tasks or execute its recommended fallback.

The model-backed examples use one deterministic two-file synthetic fixture, not issue-resolution tasks. Client evaluation separated API execution from model-selected use. A scripted consumer passed 14 response cases; 11 installed-server checks covered eight discovery/diagnostic requests. All model attempts are retained. Trial~1 stopped at missing authority. Trial~2 forwarded the authority path and passed readiness, but Codex denied its next call because client approval was unavailable. With explicit, tool-specific operator preapproval, trial~3 completed four Codex~0.153.3 turns: readiness, fresh success, identified reuse and fresh verification. All four interpretations were correct; two separate fresh container executions agreed. The fifth, model-selected request invented an unsupported argument and was rejected before execution. The model declined success, but the protocol failed and stopped; the remaining decision and six interpretation cases were not run. Thus the core lifecycle passed, not the entire trial. The sparse input-parameter description motivated a separately recorded documented-use example; it does not establish the cause of the model error. That prospectively frozen demonstration supplied one neutral API card to both decision prompts. After non-model readiness and seeding, the model selected normal reuse for a status-only need and fresh verification for a fresh-evidence need; both decisions and two independent fresh checks passed. Six fixed no-tool interpretation cases also passed, including fresh failure, refusal, and rejecting a hit when fresh diagnostics were required. All eight model turns are retained, without retries. This documents bounded use, not a causal documentation effect or population accuracy. An account-free laboratory guide additionally exercises missing-authority refusal, authorized readiness, fresh success, reuse and verification through the actual server. After an external-adapter check, a clean public clone and new external installation reproduced the published procedure with 36 matching runtime modules. The recorded installation took 18.6~s and the server check 11.3~s, excluding cloning, prior Docker setup and researcher preparation. These are researcher-executed checks, not independent user observations.

All model calls used the CLI's default model; the transcripts do not identify the backing model. These are observations of the recorded client configuration, not a model-population accuracy estimate. The operating guide separates laboratory examples from operator-approved application use.

\subsection{Real AI-workflow observations}
We also provide a read-only analysis of the public SWE-rebench OpenHands trajectories \cite{swerebenchtrajectories,badertdinov2025swerebench}. After inspecting ten separate pilot rows, a fixed SHA-256 ordering selected 128 rows from the remaining 67,064. Selection did not use repository, resolution, repetition, or performance. The raw responses, source notices, retrieval receipts, exclusions, and deterministic parser are available. An initial rate-limited attempt is retained, followed by a complete recovery of the same selected rows.

Of the 128 selected episodes, 122 had the required action/observation pairing; six were excluded for unmatched execution records, not replaced. The valid set spans 104 repositories and 120 issues. It contains 745 supported direct pytest invocations and 120 exact-command repeat pairs. All repeats have intervening barriers: 112 contain a reported editor mutation and eight contain other unmeasured effects. These are observations about recorded actions, not 120 cache hits. The parser excludes compound-shell commands from the direct-command category and preserves clipped outputs and missing exit markers.

This example makes a practical distinction visible: repeated command text is an inadequate cache key for source-editing agents. Even an apparent restoration requires independent source and environment reconstruction. A purposively selected django-environ episode provides a positive, bounded mechanism example. We reconstructed its recorded source edits, added a fresh check at the intermediate edited state, and ran four controlled requests (Table~\ref{tab:state-rejoin}). After the inverse edit, the complete declared source/runtime identity and original target's cache key rejoined; the restored request returned one explicitly identified reused success, agreeing with an independent fresh 14-test outcome. The intervening collection command selected no tests and exited 5; it remained a fresh failure, not a reusable success.

\begin{table}[htbp]
\centering\small
\begin{tabular}{llrr}
\toprule
Controlled request & Actual response & Exit & Fresh nodes \\
\midrule
Seed target & \code{MISS\_EXECUTED} & 0 & 14 \\
Edited target (extra check) & \code{MISS\_EXECUTED} & 0 & 15 \\
Restored collection & \code{MISS\_FAILED} & 5 & 0 \\
Restored target & \code{HIT\_REUSED} & 0 & 14 \\
\bottomrule
\end{tabular}
\caption{One source-restoration case. The 15-test request is additional counterfactual instrumentation, not a test call made by the recorded agent. Every row includes a separate fresh full-target oracle.}
\label{tab:state-rejoin}
\end{table}

This standalone case ran after the timing VM interruption, with independently checked source bindings. Retained failed attempts exposed a missing Git-fixture marker; a separately versioned producer initialized local metadata while preserving the baseline's masking guard. The successful run uses Python 3.12.14 and pytest 9.1.1, not the recorded Python 3.9 environment. It demonstrates this content-restoration mechanism, not historical output replay, unchanged autonomous-agent behavior, or population-wide benefit.

\section{Impact}
\subsection{Reproducible validation and measured cost}
Researchers can change a target or input boundary while preserving fresh-oracle comparisons and explicit denominators. The kit retains raw invocations, source hashes, fixed manifests, failures, setup receipts, and analysis scripts, making qualification behavior and unfavorable costs inspectable.

Codex assisted software, experiment and analysis implementation and manuscript preparation. Separately implemented validators recompute outcomes from retained raw records; input-inventory and fresh-execution oracles do not substitute cached outcomes for their reference. These are automated cross-checks, not independent human replication.

Adoption requires qualifying the result-only contract, checking agreement and failure behavior, then comparing full sequences with direct execution, including available setup cost. Fast hits alone are insufficient.

The original comparison used five pinned Python libraries: packaging, pycparser, Colorama, Click, and more-itertools. Ordinary pytest is the uncached reference; disabling only the optimized hit path gives a snapshot-first ablation, not an independent competing cache. Two reversed-order seven-request sequences per library yielded 70 fresh-result agreements and 40 optimized-arm whole-task hits. The original mean conditional hit latency reduction against the ablation was 67.3--71.0\%. Complete-sequence direct/optimized ratios were 0.85--2.09. For more-itertools, optimized execution was 16.6\% slower overall than snapshot-first reuse, despite faster hits. These results are retained, including large lifecycle delays and their order sensitivity.

To examine timing stability, a prospectively recorded extension used six independently initialized blocks on each of the four previously observed short subjects, covering every ordering of direct pytest, snapshot-first reuse, and optimized reuse. The seven request states were unchanged. All arms used the same 900-second test-execution bound and container resource limits. Docker control phases in the direct/oracle helpers used 30-second limits, whereas product inspection and cleanup used 60-second limits; complete invocation timings include these different lifecycle paths. Timings cover complete helper/API invocations; host telemetry and receipt writing are outside every timed boundary. Dependency preparation and block setup are separately retained. This is a replication on seen subjects, not a new holdout or a replacement for more-itertools.

Across 24 completed planned blocks, all 168 requests agreed with their fresh full-target checks and showed the expected cache behavior, including 96 optimized hits and 72 non-hit requests. Table~\ref{tab:replication} reports full-sequence and setup-inclusive ratios, all-block spread, and conditional median-hit savings separately. No completed block or outlying invocation within those blocks is dropped; interrupted-attempt costs are reported separately. The complete-block record contains 504 timed arm invocations and 168 separate fresh-oracle captures. The original five-subject results are not pooled with this extension. A VirtualBox host assertion interrupted execution after 21 complete blocks and part of Packaging block four. A recorded amendment reran only the unfinished preselected blocks four through six using the unchanged driver. Original failures, partial records, crash remnants and recovery preflight failures remain available. This is recovered coverage, not uninterrupted completion of the no-retry protocol. Packaging's unflushed original setup total is unavailable, so its setup-inclusive ratio is omitted. The interrupted attempt additionally retains 1 complete request, 1 partial request and an outcome-less invocation directory. Including its surviving arm costs changes Packaging's direct/optimized ratio to 1.029.

\begin{table}[htbp]
\centering\small
\begin{tabular}{lrrrr}
\toprule
Subject & $D/F$ & $D/F$ + setup & Block range & Hit saving (\%) \\
\midrule
click & 1.05 & 1.05 & 0.89--1.64 & 68.1 \\
pycparser & 0.89 & 0.91 & 0.78--1.04 & 68.8 \\
colorama & 1.01 & 1.01 & 0.98--1.11 & 67.6 \\
packaging & 1.04 & --- & 0.88--1.32 & 70.8 \\
\bottomrule
\end{tabular}
\caption{Complete replicated sequence cost versus conditional hit benefit. $D/F$ is summed direct time divided by summed optimized time; above one favors reuse. Setup-inclusive ratios allocate common dependency preparation and each block's materialization to both arms. Block range retains all six complete sequences. Hit saving compares median optimized and snapshot-first hit latency, not end-to-end agent time.}
\label{tab:replication}
\end{table}

A post-hoc operating-region analysis retains every completed block and holds within-category composition fixed. Let $D_h,F_h$ and $D_m,F_m$ denote mean complete direct/optimized costs in hit and non-hit categories. For a hypothetical hit fraction $p$, fixed category means give the equality point
\[p^*=\frac{F_m-D_m}{(F_m-D_m)+(D_h-F_h)}.\]
The crossovers for Click, pycparser, Colorama and Packaging are 54.25\%, 63.58\%, 56.30\%, 55.06\%, respectively; larger fractions favor reuse in this fixed-mixture calculation. The imposed fraction is 57.14\%, explaining why full-sequence ratios can lie near or below one despite fast hits. All six block-level crossovers per subject are archived; their ranges are descriptive, not confidence intervals. The two fingerprints account for 73.0--75.9\% of optimized hit time; their inclusive enclosing timer is not double-counted. This is not a measured agent hit frequency or a deployable admission policy. Different miss/failure mixtures and host costs change the threshold. Common setup cancels from equality; additional deployment/review costs remain unmeasured, not zero.

An independent standard-library input-inventory oracle checks boundary changes without importing ZeroRun's fingerprint implementation. Linux validation covers 32 production/oracle comparisons and six explicit oracle-scope refusals. Cases include same-size content changes with restored modification time, helper/configuration/data edits, additions, deletion, renaming, directory modes, restoration, command and environment changes, missing inputs, and symlinks. A new \emph{undeclared} top-level file changes neither declared identity: the test documents this limitation rather than claiming automatic closure discovery. These checks establish implementation behavior within the tested boundary, not universal determinism.

The source-bound engineering record includes successful Python 3.10--3.14 CI regression runs and a Windows run with 1,593 passing tests, 55 skips, and 19 passing subtests. Selected public-package tests separately passed 465 tests with 21 skips and eight passing subtests; this is not a second run of the entire development suite. A separate 100-repository CLI/skill/MCP integration exercise returned 99 passes and one safe refusal, not 100 upstream test-suite executions or autonomous AI tasks. A five-library revision evaluation recorded 200 fresh comparisons but no accepted fine-grained node reuse; the pre-existing product performance gates remain unmet. These distinct evaluations are not pooled into a single reliability percentage.

\subsection{Scope, reuse, and limitations}
For an integrator, the benefit is a small Python package with a result contract and review boundary that can be inspected before enabling reuse. For a researcher, the benefit is an executable evaluation kit linking an interface response to exact implementation bytes, fresh outcomes, and full invocation cost. The public MIT release permits adaptation and commercial reuse of ZeroRun-owned code; third-party data and software keep their own licenses.

The evidence does not establish the prevalence of reusable tasks, population-wide speedup, model-token savings, improved patch quality, user productivity, or commercial demand. The trajectory audit is descriptive and resource-bounded; repository and issue clustering is reported. Its zero barrier-free pairs do not prove zero state-restoration opportunities, but they prevent interpreting repetition as measured eligibility. The controlled library targets are convenience workloads on one constrained VM, and original and replicated observations are not statistically independent datasets. The VM shared a Windows host with publication and validation tooling; host activity outside the VM was not fully isolated. Complete-block spreads and influential observations are reported without treating thousands of nested test nodes as independent samples.

No external users, production deployments, or third-party publications using ZeroRun are claimed. The client results are controlled examples. Natural-task agent behavior, generalization across clients and models, and whether qualified tasks occur often enough to benefit complete workflows remain unmeasured.

\section{Conclusions}
ZeroRun provides reusable software for explicitly identified test-status reuse in AI coding-tool integrations. Its whole-file hit path removes execution staging when no code will run, while retaining full identity checks and the existing fresh path. The public package combines runnable examples, fresh-oracle comparisons, independent boundary checks, and a real-trajectory observation tool. The results support conditional utility and expose overhead and applicability limits. The release is a research and integration artifact, not evidence of universal AI acceleration.

\section*{Data and software availability}
Code, research data and reproduction instructions are deposited in the public ZeroRun repository \cite{kapoor2026zerorunartifact}. Table 1 pins the release; its manifest binds runtime revision \code{86f42289c2f59a74b1f642f0b20d1a26b3d55e57}. Supplementary material 1, \code{ZeroRun\_SoftwareX\_reviewer.zip}, contains the inventoried software, raw evidence and analysis tools; its README explains installation and reproduction. Public trajectory responses retain CC-BY-4.0 notices and source attribution. Private cache-authentication keys and operator authority are excluded.

\section*{Funding}
This research received no specific grant from funding agencies in the public, commercial, or not-for-profit sectors.

\section*{Declaration of competing interests}
Jishan Kapoor develops ZeroRun and has an interest in its potential commercialization. No external funding or established commercial adoption is reported.

\section*{Declaration of generative AI and AI-assisted technologies}
OpenAI Codex was used extensively to assist software and experiment implementation, analysis, literature checking, manuscript drafting, and document preparation. It is not an author. The sole human author is responsible for reviewing the evidence, verifying references and claims, and approving the submitted version.

\begingroup
\footnotesize
\bibliographystyle{elsarticle-num}
\bibliography{references}
\endgroup
\section*{Current executable software version}
ZeroRun 0.5.1, Python command-line package. The public release includes installation instructions and automated checks; Linux/amd64 Docker execution is required for the demonstrated whole-task reuse mode.
\end{document}
